\rok{Јануар 2024.}{B - Поправни тест првог теста}

Први поправни тест из Алгоритама и структура података

Први практични тест одржан 19.01.2024.

\zadatak{1}{Додавање елемената у ред (Queue Enqueue)}
Написати методу која имплементира стандардан начин за додавање елемената у ред.

\textbf{Додатне напомене за решавање задатка:}
\begin{itemize}
	\item Алгоритам написати у оквиру закоментарисане методе \texttt{enqueue(int value)}.
	\item Поменута метода налази се у оквиру класе \texttt{Queue} која се налази у придруженом template-у.
	\item Обезбедити код у Main методи за тестирање алгоритма.
\end{itemize}



\zadatak{2}{Паран број непарних вредности у колонама матрице}
Написати алгоритам који као улаз има матрицу бројева $N \times M$, број редова и број колона те матрице, и који проверава да ли се у колони елемената налази паран број непарних вредности. Алгоритам враћа низ са онолико елемената колико је колона у матрици, тако да је на $i$-том индексу овог низа јединица уколико је у $i$-тој колони матрице паран број непарних вредности, у супротном на $i$-том индексу излазног низа треба бити нула.

\textbf{Додатне напомене за решавање задатка:}
\begin{itemize}
	\item Имплементирати алгоритам.
	\item У Main методи обезбедити тестирање, позивом методе са параметрима који су дати у примерима.
	\item 0 се посматра одвојено - није ни паран, ни непаран број.
\end{itemize}

\textbf{Додатни захтеви за решавање задатка:}
\begin{itemize}
	\item Захтеван потпис методе:
	\begin{Verbatim}[fontsize=\footnotesize, numbers=left, xleftmargin=10mm]
public static int[] checkEvenOddColumns(int[,] matrix, int rows, int cols)
	\end{Verbatim}
\end{itemize}

\textbf{Примери:}
\begin{itemize}
	\item \textbf{Улаз 1:} $\begin{bmatrix} 2 & 0 & 4 & 0 \\ 3 & 5 & 7 & 9 \\ 0 & 1 & 1 & 2 \\ 1 & 2 & 3 & 4 \end{bmatrix}$, 4, 4
	\item \textbf{Излаз 1:} \texttt{[1, 1, 0, 1]}
	\item \textbf{Улаз 2:} $\begin{bmatrix} 2 & 4 & 4 \\ 1 & 4 & 3 \\ 1 & 1 & 2 \end{bmatrix}$, 3, 3
	\item \textbf{Излаз 2:} \texttt{[0, 0, 0]}
\end{itemize}



\zadatak{3}{Проналажење максимално удаљених дупликата}
Претпоставити да постоји једноструко повезана листа која садржи $2n$ елемената, од којих је $n$ различитих вредности и сваки од $n$ елемената има један дупликат у низу. Написати алгоритам који враћа удаљеност дупликата који су од свих дупликата најудаљенији. Удаљеност елемената одређена је бројем чворова који се налазе између њих.

\textbf{Додатне напомене за решавање задатка:}
\begin{itemize}
	\item Захтеван потпис методе:
	\begin{Verbatim}[fontsize=\footnotesize, numbers=left, xleftmargin=10mm]
public static int findMaximumDuplicateDistance(LinkedList list)
	\end{Verbatim}
\end{itemize}

\textbf{Примери:}
\begin{itemize}
	\item \textbf{Улаз:} \texttt{LinkedList: 1 -> 2 -> 3 -> 4 -> 2 -> 3 -> 0 -> 4 -> 0 -> 1}
	\item \textbf{Излаз:} 8
	\item \textbf{Улаз:} \texttt{LinkedList: 2 -> 2 -> 3 -> 5 -> 3 -> 5}
	\item \textbf{Излаз:} 1
	\item \textbf{Улаз:} \texttt{LinkedList: 9 -> 9 -> 3 -> 3 -> 4 -> 4}
	\item \textbf{Излаз:} 0
\end{itemize}



\sablon{Шаблон првог поправног теста јануар 2024. група B}{2023/Januar/GrupaB/AISP2023_PopravniTest1_GrupaB.linq}
